% =====================================================================
%  Sezione aggiuntiva: il caso con carico (C3P)
%  Usa i comandi già definiti nel preambolo del report:
%    \Co \Ct \Cth, \figph{...}{...}{...}{...}, \SI{}{} di siunitx.
%  \Cthp lo definisce qui sotto, così il file funziona anche da solo.
% =====================================================================

\providecommand{\Cthp}{\textsf{C3P}}

\section{Trasporto di un carico: il caso \Cthp}
\label{sec:c3p}

Tutte le misure riportate finora valutano l'assetto del corpo come una
grandezza da minimizzare, senza dire \emph{a che cosa serva} tenerlo fermo.
L'ultima configurazione provata colma questa lacuna: lo stesso controllore
\Cth{} con un carico attivo, libero di muoversi.

\subsection{Che cos'è, e perché non è un quarto controllore}

\Cthp{} è \Cth{} con il carico del modello abilitato: un cubo da
\SI{1}{\kilogram} appoggiato — non vincolato — su un vassoio da
\SI{0.1}{\kilogram} solidale al corpo. Sono \SI{1.1}{\kilogram} contro i
\SI{1.585}{\kilogram} del robot scarico: un incremento del \textbf{69\%} della
massa.

La legge di controllo è identica a quella di \Cth{}: stessi guadagni, stessa
saturazione, stesso limitatore di pendenza, nessuna ritaratura. Anche
l'impianto è lo stesso file: la differenza fra \Cth{} e \Cthp{} è lo stato
\texttt{Commented} di quattro blocchi (i due solidi, il giunto a sei gradi di
libertà e il contatto fra cubo e vassoio), commutato in memoria da una
routine dedicata prima della simulazione. Nessuna variante del modello viene
salvata su disco.

\begin{quote}
\Cthp{} non è quindi un termine di paragone per \Co{} e \Ct{}: è un
\emph{caso di carico} di \Cth{}. Si affianca a \Cth{} scarico, e alla domanda
``il controllore regge anche con il pacco?'' --- non a quella ``quale
controllore è migliore?''.
\end{quote}

\subsection{Perché il carico è libero, e perché questo cambia la misura}

Il cubo non è saldato al vassoio: è collegato da un giunto a sei gradi di
libertà e tenuto su dal solo contatto, con attrito. Può quindi scivolare,
ruotare e, in linea di principio, cadere.

Questa scelta trasforma l'assetto del corpo da numero da minimizzare a
\emph{requisito funzionale}. Un beccheggio di qualche grado non è più soltanto
una metrica peggiore: è un'inclinazione del vassoio che il carico vede, e che
può bastare a farlo camminare fino al bordo. È l'unica misura della campagna in
cui l'anello d'assetto ha un obiettivo esterno al controllo stesso, ed è anche
il motivo per cui una piattaforma a zampe viene scelta al posto di una ruotata:
non per camminare, ma per portare qualcosa dove le ruote non arrivano.

Il criterio, dichiarato prima di lanciare le run, è duplice:

\begin{enumerate}
\item il task va portato a termine --- salita riuscita per T4, dosso
attraversato per T4D, e per T6, dove in \SI{25}{\second} nessuno dei tre
controllori esce dal percorso nemmeno scarico, nessun ribaltamento, deriva
sotto \SI{0.183}{\meter} e distanza almeno il 90\% di quella di \Cth{};
\item il pacco deve restare sul vassoio: spostamento massimo dal riferimento
sotto \SI{50}{\milli\meter}, che è il gioco fra cubo e vassoio. Sotto
\SI{20}{\milli\meter} lo classifichiamo come fermo.
\end{enumerate}

\subsection{Risultati}

Le prove sono state eseguite sui tre task con terreno non piano --- rampa
(T4), dosso (T4D) e percorso a ostacoli (T6) --- cioè quelli in cui il carico
è effettivamente sollecitato.

\begin{table}[htbp]
\centering
\caption{\Cth{} scarico contro \Cthp{} carico, sugli stessi task e con la
stessa legge di controllo. Costo di trasporto ed energia sono esclusi perché
normalizzati sulla massa senza carico. Rumore di fondo non misurato su questa
configurazione: nessun verdetto di dichiarabilità.}
\label{tab:c3p}
\small
\begin{tabular}{@{}llrrl@{}}
\toprule
task & grandezza & \Cth & \Cthp & \\
\midrule
T4   & distanza [m]                    & 2.42   & 2.35   & \\
     & deriva laterale max [m]         & 0.0874 & 0.0276 & $-68\%$ \\
     & beccheggio max [deg]            & 1.27   & 1.33   & $+5\%$ \\
     & rollio max [deg]                & 0.426  & 0.441  & $+4\%$ \\
     & coppia di picco [N\,m]          & 2.70   & 4.13   & $+53\%$ \\
     & coppia RMS, giunto peggiore [N\,m] & 0.670 & 1.00  & $+49\%$ \\
     & spostamento del pacco           & ---    & \SI{19}{\milli\meter} & fermo \\
\midrule
T4D  & distanza [m]                    & 3.70   & 3.71   & \\
     & deriva laterale max [m]         & 0.131  & 0.0353 & $-73\%$ \\
     & beccheggio max [deg]            & 1.98   & 1.68   & $-15\%$ \\
     & rollio max [deg]                & 0.611  & 0.446  & $-27\%$ \\
     & coppia di picco [N\,m]          & 3.12   & 3.51   & $+13\%$ \\
     & coppia RMS, giunto peggiore [N\,m] & 0.589 & 0.956 & $+62\%$ \\
     & spostamento del pacco           & ---    & \SI{14}{\milli\meter} & fermo \\
\midrule
T6   & distanza [m]                    & 2.99   & 3.08   & \\
     & deriva laterale max [m]         & 0.0582 & 0.0154 & $-74\%$ \\
     & beccheggio max [deg]            & 4.44   & 3.53   & $-20\%$ \\
     & rollio max [deg]                & 2.78   & 2.07   & $-26\%$ \\
     & coppia di picco [N\,m]          & 4.18   & 4.62   & $+11\%$ \\
     & coppia RMS, giunto peggiore [N\,m] & 0.629 & 0.972 & $+55\%$ \\
     & spostamento del pacco           & ---    & \SI{17}{\milli\meter} & fermo \\
\bottomrule
\end{tabular}
\end{table}

\paragraph{Il carico arriva a destinazione.} Su tutti e tre i task il pacco
resta fermo sul vassoio: spostamento massimo \SI{19}{\milli\meter} (T4),
\SI{14}{\milli\meter} (T4D), \SI{17}{\milli\meter} (T6), contro un gioco
disponibile di \SI{50}{\milli\meter}. Il margine è di un fattore due e mezzo
sul caso peggiore, e il pacco non è mai arrivato al bordo. La componente
verticale dello spostamento resta sotto \SI{6}{\milli\meter}, cioè il cubo non
rimbalza: scivola di pochi millimetri e si assesta.

\paragraph{L'assetto non peggiora.} Il risultato controintuitivo è che
aggiungere il 69\% della massa non degrada l'assetto. Sulla rampa beccheggio e
rollio crescono di un $4$--$5\%$, cioè restano fermi entro quello che ci
aspettiamo dalla ripetibilità; sul dosso e sul percorso a ostacoli
\emph{migliorano}, fino al $-27\%$ di rollio su T4D.

\paragraph{La deriva laterale cala di un fattore tre o quattro.} È l'effetto
più marcato della tabella, ed è sistematico sui tre task: da $0.0874$ a
\SI{0.0276}{\meter}, da $0.131$ a \SI{0.0353}{\meter}, da $0.0582$ a
\SI{0.0154}{\meter}. Una spiegazione plausibile è che la massa aggiuntiva
aumenti la forza normale ai piedi e quindi l'attrito disponibile, riducendo lo
strisciamento in appoggio; ma \textbf{non l'abbiamo verificata}, e la
riportiamo come ipotesi e non come conclusione. L'esperimento che la
separerebbe da un semplice effetto inerziale è ripetere le run con il cubo
vincolato al vassoio: stessa massa, nessuno scivolamento relativo.

\figph{58mm}{fig:c3p-pacco}{Spostamento del pacco sul vassoio rispetto alla
posizione di riferimento, per i tre task, con la soglia di
\SI{50}{\milli\meter} corrispondente al gioco fra cubo e
vassoio.}{grafico spostamento del pacco vs tempo, tre curve T4/T4D/T6, con la soglia a 50 mm}

\subsection{Il prezzo si paga ai giunti}

Il carico non è gratis, e il conto arriva tutto sulla coppia richiesta. La
coppia efficace del giunto più caricato passa da circa \SI{0.6}{\newton\meter}
a circa \SI{1.0}{\newton\meter}, cioè dal $40$--$45\%$ al \textbf{$64$--$67\%$
della coppia di stallo} dell'AX-12A (\SI{1.5}{\newton\meter}). Rispetto al
carico continuo che il costruttore raccomanda per un moto stabile, un quinto
dello stallo, si passa da $2.2\times$ a $3.3\times$. Le punte d'urto arrivano
a $2.3$--$3.1$ volte lo stallo.

\begin{quote}
\Cthp{} è quindi il caso in cui lo scarto fra simulazione e hardware pesa di
più. Nella campagna gli attuatori sono ideali e la coppia è una reazione, non
un ingresso: quello che misuriamo è ciò che il controllore \emph{chiede}. Su
servo reali questa è la configurazione che per prima smetterebbe di
funzionare.
\end{quote}

\subsection{Limiti di questa misura}

\begin{description}[style=unboxed,leftmargin=0pt,font=\normalfont\itshape]

\item[Nessun verdetto di dichiarabilità.] Il rumore di fondo non è stato
misurato su \Cthp{}, quindi la regola del $3\times$ della
Sezione~\ref{sec:method} non è applicabile. Nessuna delle differenze della
Tabella~\ref{tab:c3p} è dichiarata: sono riportate come osservazioni. L'unico
esito binario, e quindi non soggetto al rumore, è che i tre task sono stati
portati a termine con il pacco a bordo.

\item[Energia e costo di trasporto non confrontabili.] Sono normalizzati sulla
massa del robot scarico, quindi fra \Cth{} e \Cthp{} misurerebbero soprattutto
la differenza di normalizzazione. Sono esclusi dalla tabella.

\item[Un solo carico.] Una sola massa, una sola posizione sul corpo, un solo
coefficiente d'attrito fra cubo e vassoio. Non c'è una spazzata, quindi non
sappiamo dove stia il limite: sappiamo che a $+69\%$ di massa il sistema
regge, non qual è il massimo.

\item[Solo tre task.] \Cthp{} è stato provato su rampa, dosso e percorso a
ostacoli, cioè dove il carico è sollecitato. Non in piano, in curva né sotto
impulso laterale --- e quest'ultimo, vista la lacuna comune a tutti e tre i
controllori, sarebbe il caso più severo per un pacco libero.

\item[T6 significa una cosa precisa.] Su quel task nessuno dei tre controllori
completa il percorso nei \SI{25}{\second}, nemmeno \Cth{} scarico. ``Portato a
termine'' per \Cthp{} significa quindi: nessun ribaltamento, deriva entro la
soglia di validità, e distanza non inferiore al 90\% di quella di \Cth{}
(\SI{3.08}{\meter} contro \SI{2.99}{\meter}, cioè il 103\%). Non significa che
il percorso sia stato superato.

\end{description}
